\section{Conclusion}
\label{section:conclusion}

We report on a mechanization of the proof theory of Soprillo, a fragment of SAX
covering identity, implication, tensor, and snips, in the proof assistant
Beluga. Specifically, we mechanize the translation between Soprillo and the
sequent calculus $G_3$, and present mechanized proofs of cut admissibility, cut
elimination, and the subformula property.

\subsection{Future Work}

% Move the SNAX part to the end. Saying that SNAX should follow easily from our
% mechanization is misleading. Cut-elimination (and the subformula property)
% doesn't hold in SNAX.

The most obvious avenue for future work is to extend our mechanization to the
complete SAX system, adding connectives such as $\with$, $\lor$, and
$\mathbf{1}$. This would further demonstrate the adaptability of our
representations and proof techniques. A mechanization of SAX within adjoint
logic would similarly allow for meta-theoretic guarantees regarding the
message-passing interpretation of SAX.

% Further, a mechanization of SNAX should
% follow relatively easily from our work, given that SNAX makes eligibility a
% first-class concept building directly on the proof theory we have mechanized.

Another natural step is a mechanization of the shared memory interpretation of
SAX, along with its properties including progress and termination.
\cite{DeYoung2020} presents a logical relations argument for the latter, which
are particularly amenable to mechanization. We suspect that the dynamic
semantics of the shared memory interpretation, which takes the form of multiset
rewriting rules on configurations (allocated memory cells) can be represented as
ordered lists in Beluga. While we intended to explore this path further, our
initial challenges at properly representing eligibility in Beluga stalled the
project significantly. Following this, a possible extension would be to
mechanize SNAX, which builds on SAX by making eligibility a first-class concept
in its type theory and shares some of its properties such as type progress and
preservation.
